% !TeX program = xelatex
% UTF-8; standalone XeLaTeX/XePersian transcription.
% Source: Google Drive / Notebook / GE, handwritten pages 1--2.
% images/pic01.jpg = 20261018_194117353.jpg (1ru3bJ5jpMrl4UW4xzEymw_dlAeKTKo0S)
% images/pic02.jpg = 20261018_194121455.jpg (189RXDGlVgvon-AOsn-1biv72CvdM_o3m)
% Source wording, numbers, arrows, highlighting and legible deleted text retained.
\documentclass[11pt,a4paper]{article}
\usepackage[margin=20mm]{geometry}
\usepackage{amsmath,amssymb}
\usepackage{enumitem}
\usepackage{xcolor}
\usepackage[normalem]{ulem}
\usepackage{xepersian}
\IfFontExistsTF{Amiri}{\settextfont{Amiri}}{\settextfont{DejaVu Sans}}
\setlatintextfont{DejaVu Serif}
\setlength{\parindent}{0pt}
\setlength{\parskip}{4pt}
\setlength{\emergencystretch}{3em}
\setlist[itemize]{itemsep=3pt,topsep=3pt,leftmargin=*}
\newcommand{\en}[1]{\lr{#1}}
\newcommand{\rel}{\enspace\lr{\(\longleftarrow\)}\enspace}
\newcommand{\hl}[1]{\begingroup\setlength{\fboxsep}{1pt}\colorbox{yellow!40}{\strut #1}\endgroup}
\newcommand{\sheet}[1]{%
  \section*{گاستروانتریت ـ تصویر #1}
  {\small سربرگ چاپی: \en{Date:}\quad\en{/ /}\quad\en{Subject:}
  \quad\en{Sat. Sun. Mon. Tue. Thu. Wed. Fri.}\par
  موضوع دست‌نویس: گاستروانتریت؛ شمارهٔ دست‌نویس پایین برگه: #1.\par}}

\begin{document}

\sheet{۱}
\section*{مقایسهٔ گاستروانتریت ویروسی و باکتریال}

\subsection*{دوز عفونی}
\begin{itemize}
\item ویروس \rel کم (۱۰ تا ۱۰۰ ویروس).
\item باکتری:
  \begin{itemize}
  \item بالا \rel سالمونلا، ویبریو کلرا، \en{EColi}.
  \item متوسط \rel کمپیلوباکتر ژژونی.
  \item کم \rel شیگلا.
  \end{itemize}
\end{itemize}

\subsection*{ویژگی فصلی}
\begin{itemize}
\item ویروس: در مناطق معتدل در زمستان، در گرمسیر تمام طول سال.
\item باکتری: بیشتر در تابستان و فصل باران.
\end{itemize}

\subsection*{دورهٔ کمون}
\begin{itemize}
\item ویروس: ۱ تا ۳ روز، در نوروویروس کوتاه‌تر.
\item باکتری: ۱ تا ۷ روز (کمپیلوباکتر، \en{EColi}، شیگلا، سالمونلا).

کمتر از چند ساعت (استاف اورئوس، باسیلوس سرئوس).
\end{itemize}

\subsection*{مخزن}
\begin{itemize}
\item ویروس: بیشتر انسانی.
\item باکتری: انسانی (شیگلا، سالمونلا)، حیوانی (کمپیلوباکتر، سالمونلا، \en{EColi}).
\end{itemize}
\textbf{حاشیهٔ چپ:} آب (ویبریو).

\subsection*{تب}
\begin{itemize}
\item ویروس: روتاویروس و نوروویروس تب دارند.
\item باکتری: شایع در عوامل اسهال التهابی (سالمونلا، شیگلا).
\end{itemize}

\subsection*{استفراغ}
\begin{itemize}
\item ویروس: علامت غالب / در کودکان می‌تواند تنها علامت باشد.
\item باکتری: در باکتری‌های تولیدکنندهٔ توکسین شایع است.
\end{itemize}

\subsection*{اسهال}
\begin{itemize}
\item ویروس: شایع، اغلب غیرخونی.
\item باکتری: علامت غالب / به ندرت در موارد اسهال التهابی خون می‌شود.
\end{itemize}

\subsection*{دورهٔ بهبودی}
\begin{itemize}
\item ویروس: ۱ تا ۳ روز در نوروویروس و ساپوویروس / ۲ تا ۸ روز سایر ویروس‌ها.
\item باکتری: ۱ تا ۲، باکتری تولیدکنندهٔ توکسین / ۲ تا ۸ روز سایر باکتری‌ها.
\end{itemize}

\subsection*{تشخیص}
\begin{itemize}
\item ویروس \rel \en{RO} سایر بیماری‌ها / روتاویروس و آدنوویروس
  \sout{\en{EHC}}\footnote{خوانش اختصار انگلیسی خط‌خورده نامطمئن است.}
  ایمونواسی.
\item باکتری \rel \hl{لکوسیت و خون مدفوع}.
\end{itemize}

\clearpage
\sheet{۲}

\subsection*{درمان}
\begin{itemize}
\item ویروس: حمایتی / \hl{\en{AB} و مهارکنندهٔ حرکات روده کنتراندیکاسیون دارد}.
\item باکتری: حمایتی در اکثر بیماران کافی است. / \en{AB}
  \hl{در شیگلا، ویبریو کلرا و}
  \hl{کلستریدیوم دیفیسیل}.
\end{itemize}

\end{document}
